\begin{abstract}
The word \emph{data} is among the most heavily used and least precisely defined
terms in computing. Standard references converge on an informal characterization
of data as ``raw facts,'' but none specifies what structural properties a
collection of facts must possess before it can be addressed, compared, or
processed by a machine. This ambiguity is not merely philosophical: it
propagates into system design as schema drift, type-confusion vulnerabilities,
unsound equality checks, and serialization defects, each traceable to an
unstated assumption about what index or value set a piece of data inhabits.
This paper offers a unified account of data, representation, and interpretation
across three registers. First, we ground the concept etymologically and
lexically, tracing the word \emph{datum} from its Latin root through
contemporary technical usage to expose a consistent structural implication that
informal definitions leave tacit. Second, we formalize data as a \emph{function
from an index set to a value set}, with explicit totality, functionality, and
codomain-consistency conditions, and show how this framing subsumes sequences,
tuples, and multisets as special cases. Third, we introduce the
data--information--knowledge hierarchy as the interpretive scaffolding through
which machines, which operate at the data level alone, are made to produce
outcomes meaningful to human beings, and argue that the engineering of this
transition is the central problem of data representation. We derive four
corollaries connecting violations of the formal definition's conditions to
recurring engineering failure modes---schema drift, type confusion, unsound
equality, and serialization ambiguity---and illustrate each with a worked
example. The paper concludes by distinguishing data from representation,
arguing that the choice of representation is always contingent and external to
the data itself, and that understanding this distinction is a prerequisite for
principled reasoning about encoding, type systems, and protocols.
\end{abstract}

\paragraph{Keywords.} data; representation; interpretation; formal definition;
data--information--knowledge hierarchy; type safety; schema drift;
serialization; systems design.